% Seção 8 — Overlay descritivo com mortalidade evitável.
% Enquadramento: CAUTELA DE MENSURAÇÃO, não "desertos matam".
% Disciplina: estritamente associativo. O causal mora no paper18.
% Números rastreados a 03_analysis/04_build_age_std_mortality.py e
% mortality_crude_vs_agestd_by_quadrant.csv (% src: ... abaixo).

\section{Overlay com mortalidade evitável: por que o cross-section não enxerga o
custo do isolamento}
\label{sec:mortalidade}

Se o deserto de fluxo capta isolamento que a distância não vê, a expectativa
natural é que ele apareça também na mortalidade: municípios cujos pacientes
dependem de hubs distantes deveriam acumular mais óbitos por causas que um
sistema de saúde acessível trataria a tempo. Esta seção testa essa expectativa
no plano descritivo --- e mostra que ela não se confirma no corte transversal.
O resultado não enfraquece a tese de mensuração das seções anteriores; ele
delimita o que um mapa de desertos pode e não pode ser validado contra, e por
que a pergunta sobre vidas exige o desenho causal que vive em outro lugar.

\subsection{Medida e ajuste por idade}

Usamos mortalidade evitável na definição Nolte--McKee / OCDE 2019 (causas
``amenable'' ao cuidado, com teto etário padrão de 75 anos), construída a partir
do SIM por município de residência para 2015--2022.% src: 04_build_age_std_mortality.py (lista importada de paper18 02_build_amenable_mortality.py)
A taxa bruta, porém, é uma armadilha: os desertos --- de distância ou de fluxo
--- são sistematicamente mais pobres e mais jovens, e mortalidade bruta cai com
idade média menor por pura composição demográfica, não por melhor acesso. Para
separar composição de risco, padronizamos por idade pelo método indireto (razão
de mortalidade padronizada, SMR), usando a estrutura etária municipal do Censo
2022 em faixas quinquenais 0--74 combinada com os totais populacionais anuais
como denominador, e as taxas etárias nacionais como referência.% src: census2022_age_structure.parquet; pop_municipal; nat_rate_a em 04
A padronização indireta é preferível à direta aqui porque municípios pequenos
têm contagens etárias esparsas que tornam taxas específicas por faixa instáveis.
A varredura cobre 3,1 milhões de óbitos evitáveis abaixo de 75 anos; apenas
0,01\% caem por idade ignorada.% src: log 04 — óbitos banded 3.121.049; dropados 173

\subsection{O overlay não mostra o sinal esperado}

A Tabela~\ref{tab:mort_quadrante} cruza os quadrantes da
Seção~\ref{sec:divergencia} com a taxa bruta e a taxa padronizada por idade.
O ajuste etário fecha parte da distância --- a taxa do quadrante
\emph{perto-mas-isolado} sobe de 175 para 182 por 100 mil --- mas não inverte o
ordenamento: os três quadrantes de deserto permanecem \emph{abaixo} do quadrante
conectado, cujo SMR (0,97) é o mais alto da tabela.% src: mortality_crude_vs_agestd_by_quadrant.csv
No plano contínuo o padrão é o mesmo e, se algo, contrário à intuição: a
correlação entre a taxa padronizada e o burden de fluxo F1 é $-0{,}14$, e o
coeficiente de F1 permanece negativo ($-0{,}11$ óbito por 100 mil a cada
quilômetro de burden) mesmo condicionando a renda municipal e à distância ao
hospital mais próximo.% src: regressão OLS em 04 / bloco de correlações
Lido sem cuidado, o corte transversal sugeriria que o isolamento de fluxo
\emph{protege} --- conclusão que ninguém deveria extrair, e que ilustra
exatamente o ponto desta seção.

\begin{table}[t]
\centering
\caption{Mortalidade evitável por quadrante: bruta vs.\ padronizada por idade}
\label{tab:mort_quadrante}
\small
\begin{tabular}{lcccc}
\toprule
Quadrante & Munic. & Taxa bruta & Taxa padroniz. & SMR \\
          &        & (/100 mil) & (/100 mil)     &     \\
\midrule
Perto-mas-isolado (fluxo) & 595   & 175{,}1 & 181{,}8 & 0{,}93 \\
Longe e isolado           & 778   & 157{,}2 & 169{,}1 & 0{,}87 \\
Longe-mas-conectado       & 595   & 177{,}5 & 182{,}5 & 0{,}93 \\
Conectado (referência)    & 3.524 & 191{,}6 & 189{,}3 & 0{,}97 \\
\bottomrule
\end{tabular}
\par\smallskip
\begin{minipage}{0.86\linewidth}
\footnotesize
\textit{Nota:} Mortalidade evitável Nolte--McKee / OCDE 2019, idade $<75$,
residência, 2015--2022. Padronização indireta (SMR) com estrutura etária do
Censo 2022 e taxas etárias nacionais; taxa padronizada $=$ SMR $\times$ taxa
crua nacional (195,5/100 mil). Quadrantes definidos na
Seção~\ref{sec:divergencia} (top-quartil de cada medida).
Medianas municipais. Fonte: SIM, IBGE.
\end{minipage}
\end{table}

\subsection{Por que a associação simples engana}

Dois mecanismos, ambos de mensuração e não de acesso, empurram o sinal na
direção errada. O primeiro é composicional. Mesmo depois do ajuste etário, os
municípios conectados carregam mortalidade mais alta em geral --- a mortalidade
total bruta é de 741 por 100 mil no quadrante conectado contra cerca de 688 nos
de fluxo isolado ---,% src: bloco garbage/total-mort por quadrante
reflexo de um perfil epidemiológico urbano (mais doença cardiovascular e
crônica em idade comparável) que a padronização por faixas quinquenais reduz mas
não elimina. O ajuste por idade não é ajuste por epidemiologia.

O segundo é o sub-registro. Causas mal definidas (capítulo R do CID-10)
deslocam óbitos para fora da lista evitável, e a cobertura do SIM é menor
justamente nas regiões mais remotas. Aqui, contudo, o garbage-coding é uma
explicação fraca: a fração de óbitos R96--R99 é de 2,9\% a 4,1\% nos desertos
contra 3,3\% no conectado --- diferença pequena demais para gerar o gap
observado.% src: pct_garbage_R99 por quadrante
Descartá-la importa, porque sobra o mecanismo mais incômodo: o fluxo é
\emph{endógeno à oferta e à própria sobrevivência}. Um morador que não consegue
viajar e morre sem diagnóstico oportuno não aparece como internação em outro
município nem, muitas vezes, como óbito evitável bem classificado. A geografia
da mortalidade observada já incorpora quem o sistema perdeu antes de contá-lo; o
corte transversal não tem como recuperar o contrafactual de acesso.

\subsection{O que esta seção contribui}

A lição é metodológica e vale para além do Brasil: mapas de deserto ---
medidos em quilômetros ou em fluxo --- não devem ser validados contra a
geografia bruta da mortalidade. A ausência de gradiente descritivo não é
evidência de que o isolamento não custa vidas; é evidência de que composição
demográfica, perfil epidemiológico e seleção de sobrevivência contaminam a
associação simples nos dois sentidos. Estabelecer o custo em saúde do
isolamento exige variação plausivelmente exógena na oferta e um desenho que
isole o efeito de um choque de acesso --- precisamente o terreno do desenho de
perda-de-provedor com exposição por fluxo revelado desenvolvido em trabalho
relacionado. Este relatório, por construção, para na associação e entrega a
medida; a pergunta sobre vidas fica explicitamente reservada ao desenho causal.
